\documentclass[letterpaper, 12pt]{article}

\title{Constructors \color{red}\textbf{[KEY]}}
\author{Amiel L. Iliesi}
\date{2026-03-27}

% packages
\usepackage[hidelinks]{hyperref}
\usepackage[x11names]{xcolor}
\usepackage{minted}
\usepackage{ulem}

% commands
\newcommand{\question}[1]{\textbf{#1}\par}
\newcommand{\ans}[1]{\color{Green4}\textbf{#1}\normalcolor}

% global options
\setlength{\parindent}{0pt}
\usemintedstyle{monokai}

% environments
\newenvironment{ansblock}{\color{Green4}}{\normalcolor\vspace{1em}}

\begin{document}

\maketitle

\tableofcontents

\pagebreak
\section{Dynamic Memory}

\subsection{Dynamic Element}
\question{How do you initialize a single element of type \texttt{T}?}
\begin{minted}[bgcolor=darkgray]{cpp}
T* value = new T;
\end{minted}

\question{How do you deallocate that same element?}
\begin{minted}[bgcolor=darkgray]{cpp}
delete value;
\end{minted}

\question{When do you deallocate that element?}
\begin{ansblock}
When its use is over.
\end{ansblock}

\subsection{Dynamic Array}
\question{How do you initialize an array of length \texttt{N}, of type
\texttt{T}?}

\begin{minted}[bgcolor=darkgray]{cpp}
T* values = new T[N];
\end{minted}

\question{How do you deallocate that array?}

\begin{minted}[bgcolor=darkgray]{cpp}
delete [] values;
\end{minted}

\pagebreak
\section{Constructors}

\subsection{Rule of Three}
\question{Describe the rule of three}
\begin{ansblock}
If you need any of: \uline{destructor}, \uline{copy constructor}, 
\uline{copy assignment operator}, then you most likely need all three.
\end{ansblock}

\question{\textit{[extra credit]:} Describe the rule of five, and the rule of
zero, if you know them.}

\begin{ansblock}
\textbf{Rule of Five}: Same as rule of three, except if you need move
semantics, then the \uline{move constructor} and the
\uline{move copy assignment operator} also need to be implemented.

\vspace{1em}

\textbf{Rule of Zero}: Don't use raw memory to begin with, if possible. Use
other containers, and standard library memory management practices instead. If you
do need to use raw memory, then memory ownership should be all you do.
\end{ansblock}

\subsection{Default Constructor}
\question{What is the default construtor?}
\begin{ansblock}
It's the constructor that gets called when no arguments are supplied.
\end{ansblock}

\question{What are the two ways a default constructor can become unavailable?}
\begin{ansblock}
\begin{enumerate}
    \item Alternative construtors are supplied (EG: copy constructor,
    parameterized construtor).
    \item It is explicitly removed, like so:
    \begin{minted}[bgcolor=darkgray]{cpp}
T::T() = delete;
    \end{minted}
\end{enumerate}
\end{ansblock}

\begin{minipage}{\textwidth}
\question{How can you ask the program that default constructor be
re-established?}
\begin{ansblock}
By invoking the generation of the default constructor, like so:
\begin{minted}[bgcolor=darkgray]{cpp}
T::T() = default;
\end{minted}
\end{ansblock}
\end{minipage}

\subsection{Copy Constructor}
\question{What is a copy constructor?}
\begin{ansblock}
It is the constructor invoked when the object is set to another of the same
type on declaration. It copies all members from the other object.
\end{ansblock}

\question{Given that you have a type \texttt{T}, with a dynamic member
\texttt{int x}, write a copy constructor for \texttt{T}.}
\begin{minted}[bgcolor=darkgray]{cpp}
void T::T(const T& other) {
    x = new int;
    x = other.x;
}
\end{minted}

\subsection{Copy Assignment Operator}
\question{Given a type \texttt{T}, write the prototype for its copy assignment
operator.}
\begin{minted}[bgcolor=darkgray]{cpp}
void T::operator=(const T& other);
\end{minted}

\question{What is the difference between a copy constructor, and a copy
assignment operator?}
\begin{ansblock}
The difference is \textit{when} they are expected to be called, and as a
result, what data they are expected to have. A copy construtor is always called
when the object is totally new, so it it's definitely empty. copy assignment on
the other hand can happen at any point in the object's life cycle, so it might
have data you have to clear before overwriting.
\end{ansblock}

\subsection{Parameterized Constructors}
\question{What is a parameterized constructor?}
\begin{ansblock}
It's any constructor that takes parameters. There are no restrictions to
parameters, in number or type.

\vspace{1em}

For example, a copy constructor is one example of a parameterized constructor;
its paramter is another objecct of the same class type.
\end{ansblock}

\question{If I had a type \texttt{T} with members \texttt{char c}, and
\texttt{int x}, how would you make a parameterized constructor that sets all of
\texttt{T}'s members?}

\begin{minted}[bgcolor=darkgray]{cpp}
T::T(const char& _c, const int& _x) {
    c = _c;
    x = _x;
}
\end{minted}

\pagebreak
\section{Destructors}
\question{What is a destructor? What is its purpose?}
\begin{ansblock}
Whenever you have \emph{dynamic memory}, its life-cycle can't be automatically
determined. So it comes down to the user to make that manual determination. The
destructor says what happens when it's destructed, and the \texttt{delete}
invocation says \emph{when} that happens.
\end{ansblock}

\question{When do you need a destructor?}
\begin{ansblock}
When you have dynamic memory.
\end{ansblock}

\question{How do you invoke a destructor?}
\begin{ansblock}
When the object is \textbf{static}, it is automatically invoked when it goes
\emph{out of scope}.

\vspace{1em}

When the object is \textbf{dynamic}, it is invoked when it is
\emph{deallocated}, either by single deletion (\texttt{delete}), or array
deletion (\texttt{delete[]}).
\end{ansblock}

\end{document}